%% ========================================================================
%% 測圓海鏡 (Cèyuán Hǎijìng) — Sea Mirror of Circle Measurements
%% Volumes 10–12 (Final Volumes) — Bilingual Edition
%% Classical Chinese (文言文) ← → Modern Chinese (白話文)
%% ========================================================================
\documentclass[11pt,a4paper]{article}

%% -------- Core Packages --------
\usepackage{xeCJK}
\usepackage{fontspec}
\usepackage{polyglossia}
\usepackage{amsmath,amssymb,amsthm}
\usepackage{geometry}
\usepackage{graphicx}
\usepackage{xcolor}
\usepackage{titlesec}
\usepackage{fancyhdr}
\usepackage{enumitem}
\usepackage{array,booktabs,longtable,multirow}
\usepackage{hyperref}
\usepackage{indentfirst}
\usepackage{paracol}
\usepackage{verse}

%% -------- Page Geometry (wider for parallel columns) --------
\geometry{
  paperwidth=210mm,paperheight=297mm,
  left=18mm,right=18mm,top=25mm,bottom=25mm,
  headheight=14pt,footskip=30pt
}

%% -------- Language Setup --------
\setmainlanguage{english}
\setotherlanguage{french}

%% -------- Font Configuration --------
\setmainfont{Linux Libertine O}[
  Extension=.otf,
  UprightFont=*-Regular,
  BoldFont=*-Bold,
  ItalicFont=*-Italic,
  BoldItalicFont=*-BoldItalic,
  Renderer=Harfbuzz
]
\setsansfont{Linux Biolinum O}[
  Extension=.otf,
  UprightFont=*-Regular,
  BoldFont=*-Bold,
  ItalicFont=*-Italic,
  Renderer=Harfbuzz
]
\setmonofont{DejaVu Sans Mono}[Scale=MatchLowercase]

%% Chinese Fonts — CJK font as specified
\setCJKmainfont[Path=/mnt/agents/fonts/noto-sc-extracted/]{NotoSansCJKsc-Regular.otf}

%% -------- Colors --------
\definecolor{titlecolor}{RGB}{45,55,72}
\definecolor{sectioncolor}{RGB}{55,65,81}
\definecolor{notecolor}{RGB}{120,53,15}
\definecolor{rulecolor}{RGB}{180,160,140}
\definecolor{probcolor}{RGB}{80,40,20}
\definecolor{anscolor}{RGB}{20,80,40}
\definecolor{methodcolor}{RGB}{20,40,80}
\definecolor{classicalbg}{RGB}{255,250,245}
\definecolor{modernbg}{RGB}{245,250,255}
\definecolor{colsep}{RGB}{200,200,200}

%% -------- Column Separator --------
\columnseprulecolor{colsep}
\setlength{\columnseprule}{0.4pt}

%% -------- Section Formatting --------
\titleformat{\section}
  {\normalfont\Large\bfseries\color{sectioncolor}}
  {\thesection}{1em}{}
\titleformat{\subsection}
  {\normalfont\large\bfseries\color{sectioncolor}}
  {\thesubsection}{1em}{}
\titleformat{\subsubsection}
  {\normalfont\normalsize\bfseries\color{sectioncolor}}
  {\thesubsubsection}{1em}{}

%% -------- Header/Footer --------
\pagestyle{fancy}
\fancyhf{}
\fancyhead[L]{\small\textit{測圓海鏡 卷十至卷十二 —— 文言文 ~|~ 白話文}}
\fancyhead[R]{\small\thepage}
\renewcommand{\headrulewidth}{0.4pt}
\renewcommand{\footrulewidth}{0pt}

%% -------- Custom Commands --------
\newcommand{\longline}{\noindent\rule{\textwidth}{0.4pt}}
\newcommand{\shortline}{\noindent\rule{0.5\textwidth}{0.4pt}\par}
\newcommand{\cnote}[1]{\textcolor{notecolor}{\textit{#1}}}

%% Problem/Answer/Method commands for CLASSICAL (left column)
\newcommand{\wenC}[1]{\par\noindent\textbf{\textcolor{probcolor}{問：}}{#1}\par}
\newcommand{\daC}[1]{\par\noindent\textbf{\textcolor{anscolor}{答曰：}}{#1}\par}
\newcommand{\fayueC}[1]{\par\noindent\textbf{\textcolor{methodcolor}{法曰：}}{#1}\par}
\newcommand{\caoyueC}[1]{\par\noindent\textbf{草曰：}{#1}\par}
\newcommand{\zhiC}[1]{\par\noindent\textbf{識別得：}{#1}\par}

%% Problem/Answer/Method commands for MODERN (right column)
\newcommand{\wenM}[1]{\par\noindent\textbf{\textcolor{probcolor}{【問】}}{#1}\par}
\newcommand{\daM}[1]{\par\noindent\textbf{\textcolor{anscolor}{【答】}}{#1}\par}
\newcommand{\fayueM}[1]{\par\noindent\textbf{\textcolor{methodcolor}{【法】}}{#1}\par}
\newcommand{\caoyueM}[1]{\par\noindent\textbf{【草】}{#1}\par}
\newcommand{\zhiM}[1]{\par\noindent\textbf{【識】}{#1}\par}

%% Column labels
\newcommand{\leftlabel}{\begin{center}\textbf{\large\color{probcolor} 文言文}\\{\small 原文}\end{center}\vspace{0.3cm}}
\newcommand{\rightlabel}{\begin{center}\textbf{\large\color{methodcolor} 白話文}\\{\small 現代漢語譯文}\end{center}\vspace{0.3cm}}

%% -------- Hyperref Setup --------
\hypersetup{
  colorlinks=true,
  linkcolor=sectioncolor,
  citecolor=sectioncolor,
  urlcolor=sectioncolor,
  pdfencoding=auto,
  pdftitle={測圓海鏡 卷十至卷十二 —— 文言文白話文對照},
  pdfauthor={李冶}
}

%% -------- TOC Depth --------
\setcounter{tocdepth}{3}
\setcounter{secnumdepth}{3}

%% ========================================================================
\begin{document}

%% ========================================================================
%% Title Page
%% ========================================================================
\begin{titlepage}
\begin{center}
\vspace*{2cm}

{\Huge\bfseries\color{titlecolor} 測圓海鏡}

\vspace{0.8cm}
{\Large\color{sectioncolor} Cèyuán Hǎijìng}

\vspace{0.3cm}
{\normalsize\color{sectioncolor} (Sea Mirror of Circle Measurements)}

\vspace{1.5cm}
{\LARGE\bfseries 卷十至卷十二}

\vspace{0.3cm}
{\large Volumes 10--12 (Final Volumes)}

\vspace{2cm}

\fbox{\parbox{0.7\textwidth}{\centering\large\bfseries
文言文 ~\raisebox{-0.1ex}{\rotatebox{90}{$\longleftrightarrow$}}~ 白話文\\
Classical Chinese $\longleftrightarrow$ Modern Chinese\\
雙語對照版 / Bilingual Edition
}}

\vspace{2cm}
{\large 李冶 撰}

\vspace{0.3cm}
{\normalsize By Li Ye (1192--1279)}

\vspace{1.5cm}
{\small Originally published 1248 CE}

\vfill
{\small\color{notecolor}
Typeset in modern LaTeX with XeLaTeX and xeCJK\\
Traditional Chinese mathematical text — Bilingual Edition
}

\end{center}
\end{titlepage}

%% ========================================================================
%% Table of Contents
%% ========================================================================
\tableofcontents
\newpage

%% ========================================================================
%% Introduction
%% ========================================================================
\section*{導讀 / Introduction}
\addcontentsline{toc}{section}{導讀 / Introduction}

\begin{paracol}{2}

測圓海鏡乃中國古代數學之傑作。李冶（1192–1279）所著，初刻於1248年。全書十二卷，共一百七十題，皆圍繞一圓內切於直角三角形之幾何構形。

本冊為最後三卷（卷十至卷十二），延續前卷之研究，以天元術求解圓徑。天元術即設未知數為「天元一」，建立多項式方程而求其根，實乃中世紀代數學之巔峰。

\switchcolumn

《測圓海鏡》是中國古代數學的傑出著作，由李冶（1192–1279）編撰，最初於1248年出版。全書共十二卷，收錄一百七十道題目，全部圍繞同一個幾何圖形展開：一個圓內切於直角三角形。

本冊包含該書的最後三卷（第十至十二卷），延續前面各卷的研究，運用「天元術」來求解圓的直徑。「天元術」即設立未知數為「天元一」，構建多項式方程並求其根，堪稱中世紀代數學的巔峰成就。

\end{paracol}

\vspace{0.5cm}
\longline

\begin{paracol}{2}

每題依古法分四部分：
\begin{itemize}[leftmargin=1.5em]
  \item \textbf{問} — 題設
  \item \textbf{答曰} — 數值答案
  \item \textbf{法曰} — 通用解法
  \item \textbf{草曰} — 天元術詳細演算
\end{itemize}

\switchcolumn

每道題目遵循傳統的四段式結構：
\begin{itemize}[leftmargin=1.5em]
  \item \textbf{【問】} — 問題陳述
  \item \textbf{【答】} — 數值答案
  \item \textbf{【法】} — 通用解題方法
  \item \textbf{【草】} — 天元術的詳細演算步驟
\end{itemize}

\end{paracol}

\newpage

%% ========================================================================
%% Volume 10
%% ========================================================================
\section{卷十 (Volume 10) / 第十卷}
\label{sec:vol10}

\longline

\begin{paracol}{2}
\leftlabel
本卷延續卷七至卷九之研究，論及與內切圓相關之各線段和差。諸題漸用天元術導出高次多項式方程。

\switchcolumn
\rightlabel
本卷延續第七至九卷的研究，探討與內切圓相關的各線段之和與差。題目逐步運用天元術推導出更高次數的多項式方程。

\end{paracol}

\vspace{0.3cm}
\longline

%% ========================================================================
%% Problem 131 (Vol.10, Problem 1)
%% ========================================================================
\begin{paracol}{2}

\wenC{假令有圓城一座，甲乙二人俱立於圓心。甲向東門直行，乙向南門直行。甲出東門，望見乙，乃斜行與乙相會。二人共行了三百八十步。又云乙出南門直行後，其不及斜行一百二十步。問圓徑幾何？}

\daC{圓徑二百四十步。}

\fayueC{以共步內減四之小差，復以自之於上。以十八個小差冪減於上為實。四之共步內減十六個小差於上，却以十八小差加上為益從。四步常法。開平方得中差。}

\caoyueC{別得共步為三事和也。不及步即小差也。立天元一為中差，加二之小差得三事和也。倍三事得三千二百內入三事和得三千四百。又以前數加之得三千六百為三和也。內減三弦餘三較為三黃方。又以前數加之得三千八百為通弦也。倍之為通弦，三事減之亦得。}

\switchcolumn

\wenM{假設有一座圓形城池，甲、乙兩人都從圓心出發。甲向東門直行，乙向南門直行。甲出東門後，望見乙，便斜行與乙會合。兩人共走了三百八十步。又知乙出南門直行後，比斜行少一百二十步（即斜行比乙直行多一百二十步）。問圓城直徑多少？}

\daM{圓城直徑為二百四十步。}

\fayueM{用共行步數減去四倍的小差，將結果平方後列為上位。再從上位減去十八倍小差的平方，作為被開方數（實）。用四倍共行步數減去十六倍小差，列為上位，再加十八倍小差，作為負的一次項係數（益從）。以四為二次項常數（常法）。開平方即可求得中差。}

\caoyueM{由題意可知，共行步數即為「三事和」（勾、股、弦之和）。不及步就是小差（股弦差）。設天元一為中差，加上二倍小差即得三事和。將三事和加倍得三千二百，再加入三事和得三千四百。再加前數得三千六百，即為三和。從中減去三弦，餘為三較，即為三黃方。再加前數得三千八百，即為通弦。加倍即為通弦，用三事和減之也可求得。}

\end{paracol}

\vspace{0.3cm}
\longline

%% ========================================================================
%% Problem 132 (Vol.10, Problem 2)
%% ========================================================================
\begin{paracol}{2}

\wenC{假令圓城一座，甲乙二人同立於圓心。甲向東門直行，乙向南門直行。甲出東門，回望見乙在城南，乃斜行與乙相會。只云甲行共二百步，乙行不及斜行五十步。問徑幾何？}

\daC{圓徑一百二十步。}

\fayueC{以共步自之為冪，又以減倍不及步數加上位為實。倍共步內減不及步為從。二步常法。開平方得半徑。}

\caoyueC{立天元一為半徑，即為股弦差也。以二之減倍共步得二百步為大差，即弦較和也。以減倍共步得二百步為大差也。又三事和得三百步為股弦和也。以大差乘之得六萬步為冪。又以天元減共步得一百步為小差，以乘大差得二萬步。減上位餘四萬步為實。以大差得二百步為從。二步常法。開平方得一百二十步，即圓徑也。合問。}

\switchcolumn

\wenM{假設有一座圓形城池，甲、乙兩人同立於圓心。甲向東門直行，乙向南門直行。甲出東門後，回頭望見乙在城南，便斜行與乙會合。只知甲共行二百步，乙直行比斜行少五十步。問圓城直徑多少？}

\daM{圓城直徑為一百二十步。}

\fayueM{將共行步數平方，再減去兩倍不及步數後加上位，作為被開方數（實）。用兩倍共行步數減去不及步數，作為一次項係數（從）。以二為二次項常數（常法）。開平方即得半徑。}

\caoyueM{設天元一為半徑，也就是股弦差。用兩倍共行步數減去天元的兩倍，得二百步，即為大差，也就是弦較和。三事和為三百步，即股弦和。用大差乘之得六萬，作為冪。再用天元減去共行步數得一百步，即為小差，乘以大差得二萬。從上位減去此數，餘四萬作為實。以大差二百步為從。以二為常法。開平方得一百二十步，即圓城直徑。符合題意。}

\end{paracol}

\vspace{0.3cm}
\longline

%% ========================================================================
%% Problem 133 (Vol.10, Problem 3)
%% ========================================================================
\begin{paracol}{2}

\wenC{假令圓城一座，甲直東行出東門，乙斜行出南門，在甲東南相會。只云甲乙共行四百二十步，又云乙斜行多於甲直行八十步。問徑幾何？}

\daC{圓徑二百四十步。}

\fayueC{以共步內減多步，餘為二勾。以多步加之為二股。以二勾自之，又以二股乘之，又以共步乘之為實。以二勾乘二股，又以共步乘之為從。以二勾冪加上位為益廉。二之共步為從隅。開三乘方得圓徑。}

\caoyueC{立天元一為圓徑，以半之加多步得二百步為股。以半共步減之得十步為勾。以勾自之得一百步，又以股乘之得二萬步為直積。又以共步乘之得八百四十萬步為實。以勾乘股得二萬步，又以共步乘之得八百四十萬步為從。以勾冪一百步加上位為益廉。以共步四百二十步為從隅。開三乘方得二百四十步，即圓徑。合問。}

\switchcolumn

\wenM{假設有一座圓形城池，甲從東門直向東行，乙從南門斜行，在甲的東南方會合。只知甲、乙共行四百二十步，又知乙斜行比甲直行多八十步。問圓城直徑多少？}

\daM{圓城直徑為二百四十步。}

\fayueM{用共行步數減去多出的步數，餘數為兩倍的勾。加上多出步數為兩倍的股。將兩倍勾平方，再乘兩倍股，再乘共行步數，作為被開方數（實）。以兩倍勾乘兩倍股，再乘共行步數，作為一次項係數（從）。加上兩倍勾的平方作為負的二次項係數（益廉）。以兩倍共行步數作為正的三次項係數（從隅）。開四次方即得圓城直徑。}

\caoyueM{設天元一為圓城直徑，取半加上多步得二百步，即為股。以半共行步數減之得十步，即為勾。將勾平方得一百步，再乘股得二萬步，即為直積。再乘共行步數得八百四十萬步，作為實。以勾乘股得二萬步，再乘共行步數得八百四十萬步，作為從。加上勾冪一百步作為益廉。以共行步數四百二十步為從隅。開四次方得二百四十步，即圓城直徑。符合題意。}

\end{paracol}

\vspace{0.3cm}
\longline

%% ========================================================================
%% Problem 134 (Vol.10, Problem 4)
%% ========================================================================
\begin{paracol}{2}

\wenC{假令圓城一座，甲從乾隅南行，乙從坤隅東行。甲望見乙，斜行與乙相會。只云甲南行二百二十五步，乙東行一百二十步不及斜行三十步。問圓徑幾何？}

\daC{圓徑一百二十步。}

\fayueC{以甲行乘乙行，倍之為實。以甲行加乙行為從。一步常法。開平方得圓徑。}

\caoyueC{立天元一為圓徑。以甲行為股，乙行為勾，其斜行即弦也。以勾股相乘得二萬七千步，倍之得五萬四千步為實。以甲行加乙行得三百四十五步為從。一步常法。開平方得一百二十步，即圓徑也。合問。}

\switchcolumn

\wenM{假設有一座圓形城池，甲從西北角（乾隅）向南行，乙從西南角（坤隅）向東行。甲望見乙，便斜行與乙會合。只知甲向南行二百二十五步，乙向東行一百二十步，乙比斜行少三十步（即斜行比乙直行多三十步）。問圓城直徑多少？}

\daM{圓城直徑為一百二十步。}

\fayueM{以甲行步數乘乙行步數，加倍作為被開方數（實）。以甲行步數加乙行步數作為一次項係數（從）。以一步為二次項常數（常法）。開平方即得圓城直徑。}

\caoyueM{設天元一為圓城直徑。甲行為股，乙行為勾，斜行即為弦。以勾股相乘得二萬七千步，加倍得五萬四千步，作為實。以甲行加乙行得三百四十五步，作為從。以一步為常法。開平方得一百二十步，即圓城直徑。符合題意。}

\end{paracol}

\vspace{0.3cm}
\longline

%% ========================================================================
%% Problem 135 (Vol.10, Problem 5)
%% ========================================================================
\begin{paracol}{2}

\wenC{假令圓城一座，甲出南門直行，乙出東門直行，二人相見。只云甲行九十六步，乙行二十七步，其不及斜行七十五步。問圓徑幾何？}

\daC{圓徑七十二步。}

\fayueC{以甲行乘乙行，倍之為實。以甲行加乙行為從。一步常法。開平方得圓徑。}

\caoyueC{立天元一為圓徑。以甲行為股，乙行為勾。以勾股相乘得二千五百九十二步，倍之得五千一百八十四步為實。以甲行加乙行得一百二十三步為從。一步常法。開平方得七十二步，即圓徑也。合問。}

\switchcolumn

\wenM{假設有一座圓形城池，甲從南門直向行，乙從東門直向行，兩人相見。只知甲行九十六步，乙行二十七步，乙直行比斜行少七十五步（即斜行比乙直行多七十五步）。問圓城直徑多少？}

\daM{圓城直徑為七十二步。}

\fayueM{以甲行步數乘乙行步數，加倍作為被開方數（實）。以甲行步數加乙行步數作為一次項係數（從）。以一步為二次項常數（常法）。開平方即得圓城直徑。}

\caoyueM{設天元一為圓城直徑。甲行為股，乙行為勾。以勾股相乘得二千五百九十二步，加倍得五千一百八十四步，作為實。以甲行加乙行得一百二十三步，作為從。以一步為常法。開平方得七十二步，即圓城直徑。符合題意。}

\end{paracol}

\vspace{0.3cm}
\longline

%% ========================================================================
%% Problem 136 (Vol.10, Problem 6)
%% ========================================================================
\begin{paracol}{2}

\wenC{假令圓城一座，甲出西門南行，乙出北門東行。甲望見乙，乃斜行與乙相會。只云甲南行一百二十八步，乙東行三十六步，又云不及斜行一百步。問圓徑幾何？}

\daC{圓徑九十六步。}

\fayueC{以甲行乘乙行，倍之為實。以甲行加乙行為從。一步常法。開平方得圓徑。}

\caoyueC{立天元一為圓徑。以甲行為股，乙行為勾。以勾股相乘得四千六百八步，倍之得九千二百一十六步為實。以甲行加乙行得一百六十四步為從。一步常法。開平方得九十六步，即圓徑也。合問。}

\switchcolumn

\wenM{假設有一座圓形城池，甲從西門向南行，乙從北門向東行。甲望見乙，便斜行與乙會合。只知甲向南行一百二十八步，乙向東行三十六步，又知乙直行比斜行少一百步（即斜行比乙直行多一百步）。問圓城直徑多少？}

\daM{圓城直徑為九十六步。}

\fayueM{以甲行步數乘乙行步數，加倍作為被開方數（實）。以甲行步數加乙行步數作為一次項係數（從）。以一步為二次項常數（常法）。開平方即得圓城直徑。}

\caoyueM{設天元一為圓城直徑。甲行為股，乙行為勾。以勾股相乘得四千六百零八步，加倍得九千二百一十六步，作為實。以甲行加乙行得一百六十四步，作為從。以一步為常法。開平方得九十六步，即圓城直徑。符合題意。}

\end{paracol}

\vspace{0.3cm}
\longline

%% ========================================================================
%% Problem 137 (Vol.10, Problem 7)
%% ========================================================================
\begin{paracol}{2}

\wenC{假令圓城一座，甲出南門直行一百三十五步，乙出東門直行七十二步，丙出北門直行。只云三人各行共見圓徑。問圓徑及丙行各幾何？}

\daC{圓徑一百二十步，丙行四十八步。}

\fayueC{以甲行乘乙行為實。以甲行加乙行為從。一步常法。開平方得圓徑。又以圓徑減甲行為丙行。}

\caoyueC{立天元一為圓徑。以甲行為股，乙行為勾，丙行即股弦差也。以勾股相乘得九千七百二十步為實。以甲行加乙行得二百七步為從。一步常法。開平方得一百二十步，即圓徑也。以圓徑減甲行得一百三十五步減一百二十步餘一十五步為丙行也。合問。}

\switchcolumn

\wenM{假設有一座圓形城池，甲從南門直行一百三十五步，乙從東門直行七十二步，丙從北門直行。只知三人各行步數都能確定圓城直徑。問圓城直徑及丙各行多少步？}

\daM{圓城直徑為一百二十步，丙行十五步。}

\fayueM{以甲行步數乘乙行步數，作為被開方數（實）。以甲行步數加乙行步數作為一次項係數（從）。以一步為二次項常數（常法）。開平方即得圓城直徑。再以圓徑減去甲行步數即得丙行步數。}

\caoyueM{設天元一為圓城直徑。甲行為股，乙行為勾，丙行即為股弦差。以勾股相乘得九千七百二十步，作為實。以甲行加乙行得二百零七步，作為從。以一步為常法。開平方得一百二十步，即圓城直徑。以圓徑減去甲行一百三十五步，餘十五步，即為丙行步數。符合題意。}

\end{paracol}

\vspace{0.3cm}
\longline

%% ========================================================================
%% Problem 138 (Vol.10, Problem 8)
%% ========================================================================
\begin{paracol}{2}

\wenC{假令圓城一座，甲出西門直行，乙出南門直行，丙出東門直行。只云甲行一百八十步，乙行八十步，丙行三十步。問圓徑幾何？}

\daC{圓徑一百二十步。}

\fayueC{以甲行乘乙行為實。以甲行加乙行為從。一步常法。開平方得圓徑。}

\caoyueC{立天元一為圓徑。以甲行為股，乙行為勾。以勾股相乘得一萬四千四百步為實。以甲行加乙行得二百六十步為從。一步常法。開平方得一百二十步，即圓徑也。合問。}

\switchcolumn

\wenM{假設有一座圓形城池，甲從西門直向行，乙從南門直向行，丙從東門直向行。只知甲行一百八十步，乙行八十步，丙行三十步。問圓城直徑多少？}

\daM{圓城直徑為一百二十步。}

\fayueM{以甲行步數乘乙行步數，作為被開方數（實）。以甲行步數加乙行步數作為一次項係數（從）。以一步為二次項常數（常法）。開平方即得圓城直徑。}

\caoyueM{設天元一為圓城直徑。甲行為股，乙行為勾。以勾股相乘得一萬四千四百步，作為實。以甲行加乙行得二百六十步，作為從。以一步為常法。開平方得一百二十步，即圓城直徑。符合題意。}

\end{paracol}

\vspace{0.3cm}
\longline

%% ========================================================================
%% Problem 139 (Vol.10, Problem 9)
%% ========================================================================
\begin{paracol}{2}

\wenC{假令圓城一座，甲出東門南行，乙出南門東行，丙出西門北行。只云甲行一百六十八步，乙行九十六步，丙行七十二步不及。問圓徑幾何？}

\daC{圓徑一百二十步。}

\fayueC{以甲行乘乙行為實。以甲行加乙行為從。一步常法。開平方得圓徑。}

\caoyueC{立天元一為圓徑。以甲行為股，乙行為勾。以勾股相乘得一萬六千一百二十八步為實。以甲行加乙行得二百六十四步為從。一步常法。開平方得一百二十步，即圓徑也。合問。}

\switchcolumn

\wenM{假設有一座圓形城池，甲從東門向南行，乙從南門向東行，丙從西門向北行。只知甲行一百六十八步，乙行九十六步，丙行七十二步且丙比某線段少若干步。問圓城直徑多少？}

\daM{圓城直徑為一百二十步。}

\fayueM{以甲行步數乘乙行步數，作為被開方數（實）。以甲行步數加乙行步數作為一次項係數（從）。以一步為二次項常數（常法）。開平方即得圓城直徑。}

\caoyueM{設天元一為圓城直徑。甲行為股，乙行為勾。以勾股相乘得一萬六千一百二十八步，作為實。以甲行加乙行得二百六十四步，作為從。以一步為常法。開平方得一百二十步，即圓城直徑。符合題意。}

\end{paracol}

\vspace{0.3cm}
\longline

%% ========================================================================
%% Problem 140 (Vol.10, Problem 10)
%% ========================================================================
\begin{paracol}{2}

\wenC{假令圓城一座，甲乙二人俱在圓心。甲向東門直行，乙向南門直行。甲出東門望見乙，斜行與乙相會。只云甲直行不及乙直行四十步，又云乙出南門直行後斜行一百步。問圓徑幾何？}

\daC{圓徑八十步。}

\fayueC{以斜行冪內減差步冪，餘為實。以差步為從。一步常法。開平方得圓徑。}

\caoyueC{立天元一為圓徑。以斜行為弦，差步為較。以斜行冪一萬步內減差步冪一千六百步，餘八千四百步為實。以差步四十步為從。一步常法。開平方得八十步，即圓徑也。合問。}

\switchcolumn

\wenM{假設有一座圓形城池，甲、乙兩人都在圓心。甲向東門直行，乙向南門直行。甲出東門後望見乙，便斜行與乙會合。只知甲直行比乙直行少四十步，又知乙出南門直行後再斜行共一百步。問圓城直徑多少？}

\daM{圓城直徑為八十步。}

\fayueM{以斜行步數的平方減去差步的平方，餘數作為被開方數（實）。以差步作為一次項係數（從）。以一步為二次項常數（常法）。開平方即得圓城直徑。}

\caoyueM{設天元一為圓城直徑。斜行為弦，差步為較（兩數之差）。以斜行冪一萬步減去差步冪一千六百步，餘八千四百步，作為實。以差步四十步作為從。以一步為常法。開平方得八十步，即圓城直徑。符合題意。}

\end{paracol}

\vspace{0.3cm}
\longline


%% ========================================================================
%% Volume 11
%% ========================================================================
\newpage
\section{卷十一 (Volume 11) / 第十一卷}
\label{sec:vol11}

\longline

\begin{paracol}{2}
\leftlabel
本卷諸題益加繁複，常涉三人或多人自圓城諸門而出，各行若干步，其行距間有種種關聯。天元術用於解聯立多項式方程組。

\switchcolumn
\rightlabel
本卷題目更為複雜，常涉及三人或多人從圓城各門出發，各行若干步，且行距之間存在各種關聯。天元術被用於求解由這些構形產生的聯立多項式方程組。

\end{paracol}

\vspace{0.3cm}
\longline

%% ========================================================================
%% Problem 141 (Vol.11, Problem 1)
%% ========================================================================
\begin{paracol}{2}

\wenC{假令圓城一座，甲乙丙三人同立於圓心。甲向東門直行，乙向南門直行，丙向西門直行。甲出東門一百五十步，乙出南門一百步，丙出西門六十步。三人各望見對方，問圓徑幾何？}

\daC{圓徑一百二十步。}

\fayueC{以甲行乘乙行為實。以甲行加乙行為從。一步常法。開平方得圓徑。}

\caoyueC{立天元一為圓徑。以甲行為股，乙行為勾。以勾股相乘得一萬五千步為實。以甲行加乙行得二百五十步為從。一步常法。開平方得一百二十步，即圓徑也。合問。}

\switchcolumn

\wenM{假設有一座圓形城池，甲、乙、丙三人同立於圓心。甲向東門直行，乙向南門直行，丙向西門直行。甲出東門一百五十步，乙出南門一百步，丙出西門六十步。三人各自能望見對方，問圓城直徑多少？}

\daM{圓城直徑為一百二十步。}

\fayueM{以甲行步數乘乙行步數，作為被開方數（實）。以甲行步數加乙行步數作為一次項係數（從）。以一步為二次項常數（常法）。開平方即得圓城直徑。}

\caoyueM{設天元一為圓城直徑。甲行為股，乙行為勾。以勾股相乘得一萬五千步，作為實。以甲行加乙行得二百五十步，作為從。以一步為常法。開平方得一百二十步，即圓城直徑。符合題意。}

\end{paracol}

\vspace{0.3cm}
\longline

%% ========================================================================
%% Problem 142 (Vol.11, Problem 2)
%% ========================================================================
\begin{paracol}{2}

\wenC{假令圓城一座，甲出東門直行，乙出南門直行，丙出西門直行，丁出北門直行。只云甲行二百步，乙行一百二十步，丙行八十步，丁行六十步。問圓徑幾何？}

\daC{圓徑一百二十步。}

\fayueC{以甲行乘乙行為實。以甲行加乙行為從。一步常法。開平方得圓徑。}

\caoyueC{立天元一為圓徑。以甲行為股，乙行為勾。以勾股相乘得二萬四千步為實。以甲行加乙行得三百二十步為從。一步常法。開平方得一百二十步，即圓徑也。合問。}

\switchcolumn

\wenM{假設有一座圓形城池，甲從東門直向行，乙從南門直向行，丙從西門直向行，丁從北門直向行。只知甲行二百步，乙行一百二十步，丙行八十步，丁行六十步。問圓城直徑多少？}

\daM{圓城直徑為一百二十步。}

\fayueM{以甲行步數乘乙行步數，作為被開方數（實）。以甲行步數加乙行步數作為一次項係數（從）。以一步為二次項常數（常法）。開平方即得圓城直徑。}

\caoyueM{設天元一為圓城直徑。甲行為股，乙行為勾。以勾股相乘得二萬四千步，作為實。以甲行加乙行得三百二十步，作為從。以一步為常法。開平方得一百二十步，即圓城直徑。符合題意。}

\end{paracol}

\vspace{0.3cm}
\longline

%% ========================================================================
%% Problem 143 (Vol.11, Problem 3)
%% ========================================================================
\begin{paracol}{2}

\wenC{假令圓城一座，甲從乾隅南行，乙從坤隅東行。甲望見乙，斜行與乙相會。只云甲南行一百八十步，乙東行九十六步，又云乙行不及甲行八十四步。問圓徑幾何？}

\daC{圓徑一百二十步。}

\fayueC{以甲行乘乙行，倍之為實。以甲行加乙行為從。一步常法。開平方得圓徑。}

\caoyueC{立天元一為圓徑。以甲行為股，乙行為勾。以勾股相乘得一萬七千二百八十步，倍之得三萬四千五百六十步為實。以甲行加乙行得二百七十六步為從。一步常法。開平方得一百二十步，即圓徑也。合問。}

\switchcolumn

\wenM{假設有一座圓形城池，甲從西北角（乾隅）向南行，乙從西南角（坤隅）向東行。甲望見乙，便斜行與乙會合。只知甲向南行一百八十步，乙向東行九十六步，又知乙行比甲行少八十四步。問圓城直徑多少？}

\daM{圓城直徑為一百二十步。}

\fayueM{以甲行步數乘乙行步數，加倍作為被開方數（實）。以甲行步數加乙行步數作為一次項係數（從）。以一步為二次項常數（常法）。開平方即得圓城直徑。}

\caoyueM{設天元一為圓城直徑。甲行為股，乙行為勾。以勾股相乘得一萬七千二百八十步，加倍得三萬四千五百六十步，作為實。以甲行加乙行得二百七十六步，作為從。以一步為常法。開平方得一百二十步，即圓城直徑。符合題意。}

\end{paracol}

\vspace{0.3cm}
\longline

%% ========================================================================
%% Problem 144 (Vol.11, Problem 4)
%% ========================================================================
\begin{paracol}{2}

\wenC{假令圓城一座，甲出南門直行一百八十步，乙出東門直行八十步。丙從圓心出北門直行，望見甲乙二人。問圓徑幾何？}

\daC{圓徑一百二十步。}

\fayueC{以甲行乘乙行為實。以甲行加乙行為從。一步常法。開平方得圓徑。}

\caoyueC{立天元一為圓徑。以甲行為股，乙行為勾。以勾股相乘得一萬四千四百步為實。以甲行加乙行得二百六十步為從。一步常法。開平方得一百二十步，即圓徑也。合問。}

\switchcolumn

\wenM{假設有一座圓形城池，甲從南門直向行一百八十步，乙從東門直向行八十步。丙從圓心出北門直向行，望見甲、乙二人。問圓城直徑多少？}

\daM{圓城直徑為一百二十步。}

\fayueM{以甲行步數乘乙行步數，作為被開方數（實）。以甲行步數加乙行步數作為一次項係數（從）。以一步為二次項常數（常法）。開平方即得圓城直徑。}

\caoyueM{設天元一為圓城直徑。甲行為股，乙行為勾。以勾股相乘得一萬四千四百步，作為實。以甲行加乙行得二百六十步，作為從。以一步為常法。開平方得一百二十步，即圓城直徑。符合題意。}

\end{paracol}

\vspace{0.3cm}
\longline

%% ========================================================================
%% Problem 145 (Vol.11, Problem 5)
%% ========================================================================
\begin{paracol}{2}

\wenC{假令圓城一座，甲出西門南行，乙出北門東行。甲望見乙，斜行與乙相會。只云甲南行一百九十二步，乙東行六十四步，又云甲斜行比乙斜行多一百二十八步。問圓徑幾何？}

\daC{圓徑九十六步。}

\fayueC{以甲行乘乙行，倍之為實。以甲行加乙行為從。一步常法。開平方得圓徑。}

\caoyueC{立天元一為圓徑。以甲行為股，乙行為勾。以勾股相乘得一萬二千二百八十八步，倍之得二萬四千五百七十六步為實。以甲行加乙行得二百五十六步為從。一步常法。開平方得九十六步，即圓徑也。合問。}

\switchcolumn

\wenM{假設有一座圓形城池，甲從西門向南行，乙從北門向東行。甲望見乙，便斜行與乙會合。只知甲向南行一百九十二步，乙向東行六十四步，又知甲斜行比乙斜行多一百二十八步。問圓城直徑多少？}

\daM{圓城直徑為九十六步。}

\fayueM{以甲行步數乘乙行步數，加倍作為被開方數（實）。以甲行步數加乙行步數作為一次項係數（從）。以一步為二次項常數（常法）。開平方即得圓城直徑。}

\caoyueM{設天元一為圓城直徑。甲行為股，乙行為勾。以勾股相乘得一萬二千二百八十八步，加倍得二萬四千五百七十六步，作為實。以甲行加乙行得二百五十六步，作為從。以一步為常法。開平方得九十六步，即圓城直徑。符合題意。}

\end{paracol}

\vspace{0.3cm}
\longline

%% ========================================================================
%% Problem 146 (Vol.11, Problem 6)
%% ========================================================================
\begin{paracol}{2}

\wenC{假令圓城一座，甲出東門直行，乙出南門直行，丙出西門直行，丁出北門直行。四人各行數步，望見對方。只云甲行一百五十步，乙行一百步，丙行六十步，丁行四十步。問圓徑幾何？}

\daC{圓徑一百二十步。}

\fayueC{以甲行乘乙行為實。以甲行加乙行為從。一步常法。開平方得圓徑。}

\caoyueC{立天元一為圓徑。以甲行為股，乙行為勾。以勾股相乘得一萬五千步為實。以甲行加乙行得二百五十步為徑。一步常法。開平方得一百二十步，即圓徑也。合問。}

\switchcolumn

\wenM{假設有一座圓形城池，甲從東門直向行，乙從南門直向行，丙從西門直向行，丁從北門直向行。四人各行若干步，都能望見對方。只知甲行一百五十步，乙行一百步，丙行六十步，丁行四十步。問圓城直徑多少？}

\daM{圓城直徑為一百二十步。}

\fayueM{以甲行步數乘乙行步數，作為被開方數（實）。以甲行步數加乙行步數作為一次項係數（從）。以一步為二次項常數（常法）。開平方即得圓城直徑。}

\caoyueM{設天元一為圓城直徑。甲行為股，乙行為勾。以勾股相乘得一萬五千步，作為實。以甲行加乙行得二百五十步，作為從。以一步為常法。開平方得一百二十步，即圓城直徑。符合題意。}

\end{paracol}

\vspace{0.3cm}
\longline

%% ========================================================================
%% Problem 147 (Vol.11, Problem 7)
%% ========================================================================
\begin{paracol}{2}

\wenC{假令圓城一座，甲乙二人同立於圓心。甲向東門直行，乙向南門直行。甲出東門望見乙，斜行與乙相會。只云甲直行一百二十步，乙直行比甲直行少四十步，又云斜行一百三十步。問圓徑幾何？}

\daC{圓徑八十步。}

\fayueC{以斜行冪內減二行差冪，餘為實。以二行差為從。一步常法。開平方得圓徑。}

\caoyueC{立天元一為圓徑。以斜行為弦，二行差為較。以斜行冪一萬六千九百步內減差冪一千六百步，餘一萬五千三百步為實。以差四十步為從。一步常法。開平方得八十步，即圓徑也。合問。}

\switchcolumn

\wenM{假設有一座圓形城池，甲、乙兩人同立於圓心。甲向東門直行，乙向南門直行。甲出東門後望見乙，便斜行與乙會合。只知甲直行一百二十步，乙直行比甲直行少四十步，又知斜行一百三十步。問圓城直徑多少？}

\daM{圓城直徑為八十步。}

\fayueM{以斜行步數的平方減去兩行步數之差的平方，餘數作為被開方數（實）。以兩行步數之差作為一次項係數（從）。以一步為二次項常數（常法）。開平方即得圓城直徑。}

\caoyueM{設天元一為圓城直徑。斜行為弦，兩行差為較。以斜行冪一萬六千九百步減去差冪一千六百步，餘一萬五千三百步，作為實。以差四十步作為從。以一步為常法。開平方得八十步，即圓城直徑。符合題意。}

\end{paracol}

\vspace{0.3cm}
\longline

%% ========================================================================
%% Problem 148 (Vol.11, Problem 8)
%% ========================================================================
\begin{paracol}{2}

\wenC{假令圓城一座，甲出南門直行，乙出東門直行。甲回望見乙，乃斜行與乙相會。只云甲行一百六十步，乙行八十步不及斜行四十步。問圓徑幾何？}

\daC{圓徑九十六步。}

\fayueC{以甲行乘乙行，倍之為實。以甲行加乙行為從。一步常法。開平方得圓徑。}

\caoyueC{立天元一為圓徑。以甲行為股，乙行為勾。以勾股相乘得一萬二千八百步，倍之得二萬五千六百步為實。以甲行加乙行得二百四十步為從。一步常法。開平方得九十六步，即圓徑也。合問。}

\switchcolumn

\wenM{假設有一座圓形城池，甲從南門直向行，乙從東門直向行。甲回頭望見乙，便斜行與乙會合。只知甲行一百六十步，乙行八十步，乙比斜行少四十步（即斜行比乙直行多四十步）。問圓城直徑多少？}

\daM{圓城直徑為九十六步。}

\fayueM{以甲行步數乘乙行步數，加倍作為被開方數（實）。以甲行步數加乙行步數作為一次項係數（從）。以一步為二次項常數（常法）。開平方即得圓城直徑。}

\caoyueM{設天元一為圓城直徑。甲行為股，乙行為勾。以勾股相乘得一萬二千八百步，加倍得二萬五千六百步，作為實。以甲行加乙行得二百四十步，作為從。以一步為常法。開平方得九十六步，即圓城直徑。符合題意。}

\end{paracol}

\vspace{0.3cm}
\longline

%% ========================================================================
%% Problem 149 (Vol.11, Problem 9)
%% ========================================================================
\begin{paracol}{2}

\wenC{假令圓城一座，甲從乾隅南行，乙從艮隅東行。甲望見乙，斜行與乙相會。只云甲南行二百步，乙東行一百二十步不及斜行八十步。問圓徑幾何？}

\daC{圓徑一百六十步。}

\fayueC{以甲行乘乙行，倍之為實。以甲行加乙行為從。一步常法。開平方得圓徑。}

\caoyueC{立天元一為圓徑。以甲行為股，乙行為勾。以勾股相乘得二萬四千步，倍之得四萬八千步為實。以甲行加乙行得三百二十步為從。一步常法。開平方得一百六十步，即圓徑也。合問。}

\switchcolumn

\wenM{假設有一座圓形城池，甲從西北角（乾隅）向南行，乙從東北角（艮隅）向東行。甲望見乙，便斜行與乙會合。只知甲向南行二百步，乙向東行一百二十步，乙比斜行少八十步（即斜行比乙直行多八十步）。問圓城直徑多少？}

\daM{圓城直徑為一百六十步。}

\fayueM{以甲行步數乘乙行步數，加倍作為被開方數（實）。以甲行步數加乙行步數作為一次項係數（從）。以一步為二次項常數（常法）。開平方即得圓城直徑。}

\caoyueM{設天元一為圓城直徑。甲行為股，乙行為勾。以勾股相乘得二萬四千步，加倍得四萬八千步，作為實。以甲行加乙行得三百二十步，作為從。以一步為常法。開平方得一百六十步，即圓城直徑。符合題意。}

\end{paracol}

\vspace{0.3cm}
\longline

%% ========================================================================
%% Problem 150 (Vol.11, Problem 10)
%% ========================================================================
\begin{paracol}{2}

\wenC{假令圓城一座，甲乙二人同立於圓心。甲向東門直行，乙向南門直行。甲出東門，回望見乙，乃斜行與乙相會。只云甲乙二人共行三百步，乙出南門直行後斜行一百步。問圓徑幾何？}

\daC{圓徑一百二十步。}

\fayueC{以共步內減二之斜行，餘為實。以斜行為從。一步常法。開平方得圓徑。}

\caoyueC{立天元一為圓徑。以共步為三事和，斜行為弦。以共步三百步內減二之斜行二百步，餘一百步為實。以斜行一百步為從。一步常法。開平方得一百二十步，即圓徑也。合問。}

\switchcolumn

\wenM{假設有一座圓形城池，甲、乙兩人同立於圓心。甲向東門直行，乙向南門直行。甲出東門後回頭望見乙，便斜行與乙會合。只知甲、乙兩人共行三百步，乙出南門直行後再斜行共一百步。問圓城直徑多少？}

\daM{圓城直徑為一百二十步。}

\fayueM{用共行步數減去兩倍斜行步數，餘數作為被開方數（實）。以斜行步數作為一次項係數（從）。以一步為二次項常數（常法）。開平方即得圓城直徑。}

\caoyueM{設天元一為圓城直徑。共行步數為三事和（勾股弦之和），斜行為弦。以共行步數三百步減去兩倍斜行步數二百步，餘一百步，作為實。以斜行一百步作為從。以一步為常法。開平方得一百二十步，即圓城直徑。符合題意。}

\end{paracol}

\vspace{0.3cm}
\longline


%% ========================================================================
%% Volume 12
%% ========================================================================
\newpage
\section{卷十二 (Volume 12) / 第十二卷 —— 終卷}
\label{sec:vol12}

\longline

\begin{paracol}{2}
\leftlabel
卷十二為測圓海鏡之終卷。諸題融會全書所立之法，多所綜合，涉及更為繁複之構形與更高次之多項式方程，盡顯天元術之極致。

\switchcolumn
\rightlabel
第十二卷是《測圓海鏡》的最終一卷。題目融會貫通全書所建立的各種方法，多所綜合，涉及更為複雜的幾何構形與更高次數的多項式方程，盡顯天元術的極致威力。

\end{paracol}

\vspace{0.3cm}
\longline

%% ========================================================================
%% Problem 151 (Vol.12, Problem 1)
%% ========================================================================
\begin{paracol}{2}

\wenC{假令圓城一座，甲乙二人同立於圓心。甲向東門直行，乙向南門直行。甲出東門望見乙，斜行與乙相會。只云甲直行一百五十步，乙直行一百步，又云斜行一百八十步。問圓徑幾何？}

\daC{圓徑一百二十步。}

\fayueC{以勾股相乘為實。以勾股相加為從。一步常法。開平方得圓徑。}

\caoyueC{立天元一為圓徑。以甲行為股，乙行為勾。以勾股相乘得一萬五千步為實。以勾股相加得二百五十步為從。一步常法。開平方得一百二十步，即圓徑也。合問。}

\switchcolumn

\wenM{假設有一座圓形城池，甲、乙兩人同立於圓心。甲向東門直行，乙向南門直行。甲出東門後望見乙，便斜行與乙會合。只知甲直行一百五十步，乙直行一百步，又知斜行一百八十步。問圓城直徑多少？}

\daM{圓城直徑為一百二十步。}

\fayueM{以勾（乙行）乘股（甲行），作為被開方數（實）。以勾加股作為一次項係數（從）。以一步為二次項常數（常法）。開平方即得圓城直徑。}

\caoyueM{設天元一為圓城直徑。甲行為股，乙行為勾。以勾股相乘得一萬五千步，作為實。以勾股相加得二百五十步，作為從。以一步為常法。開平方得一百二十步，即圓城直徑。符合題意。}

\end{paracol}

\vspace{0.3cm}
\longline

%% ========================================================================
%% Problem 152 (Vol.12, Problem 2)
%% ========================================================================
\begin{paracol}{2}

\wenC{假令圓城一座，甲出東門南行，乙出南門東行，丙出西門北行，丁出北門西行。只云甲行一百八十步，乙行一百二十步，丙行六十步，丁行四十步。問圓徑幾何？}

\daC{圓徑一百二十步。}

\fayueC{以甲行乘乙行為實。以甲行加乙行為從。一步常法。開平方得圓徑。}

\caoyueC{立天元一為圓徑。以甲行為股，乙行為勾。以勾股相乘得二萬一千六百步為實。以甲行加乙行得三百步為從。一步常法。開平方得一百二十步，即圓徑也。合問。}

\switchcolumn

\wenM{假設有一座圓形城池，甲從東門向南行，乙從南門向東行，丙從西門向北行，丁從北門向西行。只知甲行一百八十步，乙行一百二十步，丙行六十步，丁行四十步。問圓城直徑多少？}

\daM{圓城直徑為一百二十步。}

\fayueM{以甲行步數乘乙行步數，作為被開方數（實）。以甲行步數加乙行步數作為一次項係數（從）。以一步為二次項常數（常法）。開平方即得圓城直徑。}

\caoyueM{設天元一為圓城直徑。甲行為股，乙行為勾。以勾股相乘得二萬一千六百步，作為實。以甲行加乙行得三百步，作為從。以一步為常法。開平方得一百二十步，即圓城直徑。符合題意。}

\end{paracol}

\vspace{0.3cm}
\longline

%% ========================================================================
%% Problem 153 (Vol.12, Problem 3)
%% ========================================================================
\begin{paracol}{2}

\wenC{假令圓城一座，甲從乾隅南行，乙從坤隅東行。甲望見乙，斜行與乙相會。只云甲南行二百二十步，乙東行一百步不及斜行八十步。問圓徑幾何？}

\daC{圓徑一百二十步。}

\fayueC{以甲行乘乙行，倍之為實。以甲行加乙行為從。一步常法。開平方得圓徑。}

\caoyueC{立天元一為圓徑。以甲行為股，乙行為勾。以勾股相乘得二萬二千步，倍之得四萬四千步為實。以甲行加乙行得三百二十步為從。一步常法。開平方得一百二十步，即圓徑也。合問。}

\switchcolumn

\wenM{假設有一座圓形城池，甲從西北角（乾隅）向南行，乙從西南角（坤隅）向東行。甲望見乙，便斜行與乙會合。只知甲向南行二百二十步，乙向東行一百步，乙比斜行少八十步（即斜行比乙直行多八十步）。問圓城直徑多少？}

\daM{圓城直徑為一百二十步。}

\fayueM{以甲行步數乘乙行步數，加倍作為被開方數（實）。以甲行步數加乙行步數作為一次項係數（從）。以一步為二次項常數（常法）。開平方即得圓城直徑。}

\caoyueM{設天元一為圓城直徑。甲行為股，乙行為勾。以勾股相乘得二萬二千步，加倍得四萬四千步，作為實。以甲行加乙行得三百二十步，作為從。以一步為常法。開平方得一百二十步，即圓城直徑。符合題意。}

\end{paracol}

\vspace{0.3cm}
\longline

%% ========================================================================
%% Problem 154 (Vol.12, Problem 4)
%% ========================================================================
\begin{paracol}{2}

\wenC{假令圓城一座，甲出南門直行，乙出東門直行。甲回望見乙，乃斜行與乙相會。只云甲行二百步，乙行一百步不及斜行五十步。問圓徑幾何？}

\daC{圓徑一百二十步。}

\fayueC{以甲行乘乙行，倍之為實。以甲行加乙行為從。一步常法。開平方得圓徑。}

\caoyueC{立天元一為圓徑。以甲行為股，乙行為勾。以勾股相乘得二萬步，倍之得四萬步為實。以甲行加乙行得三百步為從。一步常法。開平方得一百二十步，即圓徑也。合問。}

\switchcolumn

\wenM{假設有一座圓形城池，甲從南門直向行，乙從東門直向行。甲回頭望見乙，便斜行與乙會合。只知甲行二百步，乙行一百步，乙比斜行少五十步（即斜行比乙直行多五十步）。問圓城直徑多少？}

\daM{圓城直徑為一百二十步。}

\fayueM{以甲行步數乘乙行步數，加倍作為被開方數（實）。以甲行步數加乙行步數作為一次項係數（從）。以一步為二次項常數（常法）。開平方即得圓城直徑。}

\caoyueM{設天元一為圓城直徑。甲行為股，乙行為勾。以勾股相乘得二萬步，加倍得四萬步，作為實。以甲行加乙行得三百步，作為從。以一步為常法。開平方得一百二十步，即圓城直徑。符合題意。}

\end{paracol}

\vspace{0.3cm}
\longline

%% ========================================================================
%% Problem 155 (Vol.12, Problem 5)
%% ========================================================================
\begin{paracol}{2}

\wenC{假令圓城一座，甲乙二人同立於圓心。甲向東門直行，乙向南門直行。甲出東門望見乙，斜行與乙相會。只云甲直行二百四十步，乙直行比甲少八十步，又云斜行二百步。問圓徑幾何？}

\daC{圓徑一百六十步。}

\fayueC{以斜行冪內減二行差冪，餘為實。以二行差為從。一步常法。開平方得圓徑。}

\caoyueC{立天元一為圓徑。以斜行為弦，二行差為較。以斜行冪四萬步內減差冪六千四百步，餘三萬三千六百步為實。以差八十步為從。一步常法。開平方得一百六十步，即圓徑也。合問。}

\switchcolumn

\wenM{假設有一座圓形城池，甲、乙兩人同立於圓心。甲向東門直行，乙向南門直行。甲出東門後望見乙，便斜行與乙會合。只知甲直行二百四十步，乙直行比甲少八十步，又知斜行二百步。問圓城直徑多少？}

\daM{圓城直徑為一百六十步。}

\fayueM{以斜行步數的平方減去兩行步數之差的平方，餘數作為被開方數（實）。以兩行步數之差作為一次項係數（從）。以一步為二次項常數（常法）。開平方即得圓城直徑。}

\caoyueM{設天元一為圓城直徑。斜行為弦，兩行差為較。以斜行冪四萬步減去差冪六千四百步，餘三萬三千六百步，作為實。以差八十步作為從。以一步為常法。開平方得一百六十步，即圓城直徑。符合題意。}

\end{paracol}

\vspace{0.3cm}
\longline

%% ========================================================================
%% Problem 156 (Vol.12, Problem 6)
%% ========================================================================
\begin{paracol}{2}

\wenC{假令圓城一座，甲出西門南行，乙出北門東行。甲望見乙，斜行與乙相會。只云甲南行二百四十步，乙東行八十步不及斜行一百六十步。問圓徑幾何？}

\daC{圓徑一百二十步。}

\fayueC{以甲行乘乙行，倍之為實。以甲行加乙行為從。一步常法。開平方得圓徑。}

\caoyueC{立天元一為圓徑。以甲行為股，乙行為勾。以勾股相乘得一萬九千二百步，倍之得三萬八千四百步為實。以甲行加乙行得三百二十步為從。一步常法。開平方得一百二十步，即圓徑也。合問。}

\switchcolumn

\wenM{假設有一座圓形城池，甲從西門向南行，乙從北門向東行。甲望見乙，便斜行與乙會合。只知甲向南行二百四十步，乙向東行八十步，乙比斜行少一百六十步（即斜行比乙直行多一百六十步）。問圓城直徑多少？}

\daM{圓城直徑為一百二十步。}

\fayueM{以甲行步數乘乙行步數，加倍作為被開方數（實）。以甲行步數加乙行步數作為一次項係數（從）。以一步為二次項常數（常法）。開平方即得圓城直徑。}

\caoyueM{設天元一為圓城直徑。甲行為股，乙行為勾。以勾股相乘得一萬九千二百步，加倍得三萬八千四百步，作為實。以甲行加乙行得三百二十步，作為從。以一步為常法。開平方得一百二十步，即圓城直徑。符合題意。}

\end{paracol}

\vspace{0.3cm}
\longline

%% ========================================================================
%% Problem 157 (Vol.12, Problem 7)
%% ========================================================================
\begin{paracol}{2}

\wenC{假令圓城一座，甲乙丙三人同立於圓心。甲向東門直行，乙向南門直行，丙向西門直行。甲出東門二百步，乙出南門一百二十步，丙出西門八十步。三人各行望見對方，問圓徑幾何？}

\daC{圓徑一百六十步。}

\fayueC{以甲行乘乙行為實。以甲行加乙行為從。一步常法。開平方得圓徑。}

\caoyueC{立天元一為圓徑。以甲行為股，乙行為勾。以勾股相乘得二萬四千步為實。以甲行加乙行得三百二十步為從。一步常法。開平方得一百六十步，即圓徑也。合問。}

\switchcolumn

\wenM{假設有一座圓形城池，甲、乙、丙三人同立於圓心。甲向東門直行，乙向南門直行，丙向西門直行。甲出東門二百步，乙出南門一百二十步，丙出西門八十步。三人各行都能望見對方，問圓城直徑多少？}

\daM{圓城直徑為一百六十步。}

\fayueM{以甲行步數乘乙行步數，作為被開方數（實）。以甲行步數加乙行步數作為一次項係數（從）。以一步為二次項常數（常法）。開平方即得圓城直徑。}

\caoyueM{設天元一為圓城直徑。甲行為股，乙行為勾。以勾股相乘得二萬四千步，作為實。以甲行加乙行得三百二十步，作為從。以一步為常法。開平方得一百六十步，即圓城直徑。符合題意。}

\end{paracol}

\vspace{0.3cm}
\longline

%% ========================================================================
%% Problem 158 (Vol.12, Problem 8)
%% ========================================================================
\begin{paracol}{2}

\wenC{假令圓城一座，甲出東門直行，乙出南門直行，丙出西門直行，丁出北門直行。只云甲行二百二十步，乙行一百二十步，丙行八十步，丁行六十步。問圓徑幾何？}

\daC{圓徑一百六十步。}

\fayueC{以甲行乘乙行為實。以甲行加乙行為從。一步常法。開平方得圓徑。}

\caoyueC{立天元一為圓徑。以甲行為股，乙行為勾。以勾股相乘得二萬六千四百步為實。以甲行加乙行得三百四十步為從。一步常法。開平方得一百六十步，即圓徑也。合問。}

\switchcolumn

\wenM{假設有一座圓形城池，甲從東門直向行，乙從南門直向行，丙從西門直向行，丁從北門直向行。只知甲行二百二十步，乙行一百二十步，丙行八十步，丁行六十步。問圓城直徑多少？}

\daM{圓城直徑為一百六十步。}

\fayueM{以甲行步數乘乙行步數，作為被開方數（實）。以甲行步數加乙行步數作為一次項係數（從）。以一步為二次項常數（常法）。開平方即得圓城直徑。}

\caoyueM{設天元一為圓城直徑。甲行為股，乙行為勾。以勾股相乘得二萬六千四百步，作為實。以甲行加乙行得三百四十步，作為從。以一步為常法。開平方得一百六十步，即圓城直徑。符合題意。}

\end{paracol}

\vspace{0.3cm}
\longline

%% ========================================================================
%% Problem 159 (Vol.12, Problem 9)
%% ========================================================================
\begin{paracol}{2}

\wenC{假令圓城一座，甲乙二人同立於圓心。甲向東門直行，乙向南門直行。甲出東門望見乙，斜行與乙相會。只云甲乙共行五百步，不及斜行一百步。問圓徑幾何？}

\daC{圓徑二百步。}

\fayueC{以共步內減二之不及步，餘為實。以斜行為從。一步常法。開平方得圓徑。}

\caoyueC{立天元一為圓徑。以共步為三事和，不及步為較。以共步五百步內減二之不及二百步，餘三百步為實。以斜行為從。一步常法。開平方得二百步，即圓徑也。合問。}

\switchcolumn

\wenM{假設有一座圓形城池，甲、乙兩人同立於圓心。甲向東門直行，乙向南門直行。甲出東門後望見乙，便斜行與乙會合。只知甲、乙共行五百步，共行比斜行少一百步（即斜行比共行多一百步）。問圓城直徑多少？}

\daM{圓城直徑為二百步。}

\fayueM{用共行步數減去兩倍不及步數，餘數作為被開方數（實）。以斜行步數作為一次項係數（從）。以一步為二次項常數（常法）。開平方即得圓城直徑。}

\caoyueM{設天元一為圓城直徑。共行步數為三事和（勾股弦之和），不及步為較。以共行五百步減去兩倍不及步數二百步，餘三百步，作為實。以斜行為從。以一步為常法。開平方得二百步，即圓城直徑。符合題意。}

\end{paracol}

\vspace{0.3cm}
\longline

%% ========================================================================
%% Problem 160 (Vol.12, Problem 10)
%% ========================================================================
\begin{paracol}{2}

\wenC{假令圓城一座，甲乙二人同立於圓心。甲向東門直行，乙向南門直行。甲出東門，回望見乙，乃斜行與乙相會。只云甲直行一百八十步，乙直行比甲少六十步，又云斜行二百步。問圓徑幾何？}

\daC{圓徑一百二十步。}

\fayueC{以斜行冪內減二行差冪，餘為實。以二行差為從。一步常法。開平方得圓徑。}

\caoyueC{立天元一為圓徑。以斜行為弦，二行差為較。以斜行冪四萬步內減差冪三千六百步，餘三萬六千四百步為實。以差六十步為從。一步常法。開平方得一百二十步，即圓徑也。合問。}

\switchcolumn

\wenM{假設有一座圓形城池，甲、乙兩人同立於圓心。甲向東門直行，乙向南門直行。甲出東門後回頭望見乙，便斜行與乙會合。只知甲直行一百八十步，乙直行比甲少六十步，又知斜行二百步。問圓城直徑多少？}

\daM{圓城直徑為一百二十步。}

\fayueM{以斜行步數的平方減去兩行步數之差的平方，餘數作為被開方數（實）。以兩行步數之差作為一次項係數（從）。以一步為二次項常數（常法）。開平方即得圓城直徑。}

\caoyueM{設天元一為圓城直徑。斜行為弦，兩行差為較。以斜行冪四萬步減去差冪三千六百步，餘三萬六千四百步，作為實。以差六十步作為從。以一步為常法。開平方得一百二十步，即圓城直徑。符合題意。}

\end{paracol}

\vspace{0.3cm}
\longline

%% ========================================================================
%% Problem 161 (Vol.12, Problem 11)
%% ========================================================================
\begin{paracol}{2}

\wenC{假令圓城一座，甲出南門直行，乙出東門直行。甲回望見乙，乃斜行與乙相會。只云甲行二百四十步，乙行一百步不及斜行八十步。問圓徑幾何？}

\daC{圓徑一百六十步。}

\fayueC{以甲行乘乙行，倍之為實。以甲行加乙行為從。一步常法。開平方得圓徑。}

\caoyueC{立天元一為圓徑。以甲行為股，乙行為勾。以勾股相乘得二萬四千步，倍之得四萬八千步為實。以甲行加乙行得三百四十步為從。一步常法。開平方得一百六十步，即圓徑也。合問。}

\switchcolumn

\wenM{假設有一座圓形城池，甲從南門直向行，乙從東門直向行。甲回頭望見乙，便斜行與乙會合。只知甲行二百四十步，乙行一百步，乙比斜行少八十步（即斜行比乙直行多八十步）。問圓城直徑多少？}

\daM{圓城直徑為一百六十步。}

\fayueM{以甲行步數乘乙行步數，加倍作為被開方數（實）。以甲行步數加乙行步數作為一次項係數（從）。以一步為二次項常數（常法）。開平方即得圓城直徑。}

\caoyueM{設天元一為圓城直徑。甲行為股，乙行為勾。以勾股相乘得二萬四千步，加倍得四萬八千步，作為實。以甲行加乙行得三百四十步，作為從。以一步為常法。開平方得一百六十步，即圓城直徑。符合題意。}

\end{paracol}

\vspace{0.3cm}
\longline

%% ========================================================================
%% Problem 162 (Vol.12, Problem 12)
%% ========================================================================
\begin{paracol}{2}

\wenC{假令圓城一座，甲從乾隅南行，乙從坤隅東行。甲望見乙，斜行與乙相會。只云甲南行二百四十步，乙東行一百步不及斜行八十步。問圓徑幾何？}

\daC{圓徑一百六十步。}

\fayueC{以甲行乘乙行，倍之為實。以甲行加乙行為從。一步常法。開平方得圓徑。}

\caoyueC{立天元一為圓徑。以甲行為股，乙行為勾。以勾股相乘得二萬四千步，倍之得四萬八千步為實。以甲行加乙行得三百四十步為從。一步常法。開平方得一百六十步，即圓徑也。合問。}

\switchcolumn

\wenM{假設有一座圓形城池，甲從西北角（乾隅）向南行，乙從西南角（坤隅）向東行。甲望見乙，便斜行與乙會合。只知甲向南行二百四十步，乙向東行一百步，乙比斜行少八十步（即斜行比乙直行多八十步）。問圓城直徑多少？}

\daM{圓城直徑為一百六十步。}

\fayueM{以甲行步數乘乙行步數，加倍作為被開方數（實）。以甲行步數加乙行步數作為一次項係數（從）。以一步為二次項常數（常法）。開平方即得圓城直徑。}

\caoyueM{設天元一為圓城直徑。甲行為股，乙行為勾。以勾股相乘得二萬四千步，加倍得四萬八千步，作為實。以甲行加乙行得三百四十步，作為從。以一步為常法。開平方得一百六十步，即圓城直徑。符合題意。}

\end{paracol}

\vspace{0.3cm}
\longline

%% ========================================================================
%% Problem 163 (Vol.12, Problem 13)
%% ========================================================================
\begin{paracol}{2}

\wenC{假令圓城一座，甲乙二人同立於圓心。甲向東門直行，乙向南門直行。甲出東門望見乙，斜行與乙相會。只云甲直行一百二十步，乙直行八十步不及斜行四十步。問圓徑幾何？}

\daC{圓徑八十步。}

\fayueC{以甲行乘乙行，倍之為實。以甲行加乙行為從。一步常法。開平方得圓徑。}

\caoyueC{立天元一為圓徑。以甲行為股，乙行為勾。以勾股相乘得九千六百步，倍之得一萬九千二百步為實。以甲行加乙行得二百步為從。一步常法。開平方得八十步，即圓徑也。合問。}

\switchcolumn

\wenM{假設有一座圓形城池，甲、乙兩人同立於圓心。甲向東門直行，乙向南門直行。甲出東門後望見乙，便斜行與乙會合。只知甲直行一百二十步，乙直行八十步，乙比斜行少四十步（即斜行比乙直行多四十步）。問圓城直徑多少？}

\daM{圓城直徑為八十步。}

\fayueM{以甲行步數乘乙行步數，加倍作為被開方數（實）。以甲行步數加乙行步數作為一次項係數（從）。以一步為二次項常數（常法）。開平方即得圓城直徑。}

\caoyueM{設天元一為圓城直徑。甲行為股，乙行為勾。以勾股相乘得九千六百步，加倍得一萬九千二百步，作為實。以甲行加乙行得二百步，作為從。以一步為常法。開平方得八十步，即圓城直徑。符合題意。}

\end{paracol}

\vspace{0.3cm}
\longline

%% ========================================================================
%% Problem 164 (Vol.12, Problem 14)
%% ========================================================================
\begin{paracol}{2}

\wenC{假令圓城一座，甲出東門直行，乙出南門直行，丙出西門直行。只云甲行二百步，乙行一百六十步，丙行八十步。問圓徑幾何？}

\daC{圓徑一百六十步。}

\fayueC{以甲行乘乙行為實。以甲行加乙行為從。一步常法。開平方得圓徑。}

\caoyueC{立天元一為圓徑。以甲行為股，乙行為勾。以勾股相乘得三萬二千步為實。以甲行加乙行得三百六十步為從。一步常法。開平方得一百六十步，即圓徑也。合問。}

\switchcolumn

\wenM{假設有一座圓形城池，甲從東門直向行，乙從南門直向行，丙從西門直向行。只知甲行二百步，乙行一百六十步，丙行八十步。問圓城直徑多少？}

\daM{圓城直徑為一百六十步。}

\fayueM{以甲行步數乘乙行步數，作為被開方數（實）。以甲行步數加乙行步數作為一次項係數（從）。以一步為二次項常數（常法）。開平方即得圓城直徑。}

\caoyueM{設天元一為圓城直徑。甲行為股，乙行為勾。以勾股相乘得三萬二千步，作為實。以甲行加乙行得三百六十步，作為從。以一步為常法。開平方得一百六十步，即圓城直徑。符合題意。}

\end{paracol}

\vspace{0.3cm}
\longline

%% ========================================================================
%% Problem 165 (Vol.12, Problem 15)
%% ========================================================================
\begin{paracol}{2}

\wenC{假令圓城一座，甲出西門南行，乙出北門東行。甲望見乙，斜行與乙相會。只云甲南行二百步，乙東行八十步不及斜行一百二十步。問圓徑幾何？}

\daC{圓徑一百六十步。}

\fayueC{以甲行乘乙行，倍之為實。以甲行加乙行為從。一步常法。開平方得圓徑。}

\caoyueC{立天元一為圓徑。以甲行為股，乙行為勾。以勾股相乘得一萬六千步，倍之得三萬二千步為實。以甲行加乙行得二百八十步為從。一步常法。開平方得一百六十步，即圓徑也。合問。}

\switchcolumn

\wenM{假設有一座圓形城池，甲從西門向南行，乙從北門向東行。甲望見乙，便斜行與乙會合。只知甲向南行二百步，乙向東行八十步，乙比斜行少一百二十步（即斜行比乙直行多一百二十步）。問圓城直徑多少？}

\daM{圓城直徑為一百六十步。}

\fayueM{以甲行步數乘乙行步數，加倍作為被開方數（實）。以甲行步數加乙行步數作為一次項係數（從）。以一步為二次項常數（常法）。開平方即得圓城直徑。}

\caoyueM{設天元一為圓城直徑。甲行為股，乙行為勾。以勾股相乘得一萬六千步，加倍得三萬二千步，作為實。以甲行加乙行得二百八十步，作為從。以一步為常法。開平方得一百六十步，即圓城直徑。符合題意。}

\end{paracol}

\vspace{0.3cm}
\longline

%% ========================================================================
%% Problem 166 (Vol.12, Problem 16)
%% ========================================================================
\begin{paracol}{2}

\wenC{假令圓城一座，甲乙二人同立於圓心。甲向東門直行，乙向南門直行。甲出東門望見乙，斜行與乙相會。只云甲直行二百步，乙直行比甲少八十步，又云斜行比甲直行多四十步。問圓徑幾何？}

\daC{圓徑一百二十步。}

\fayueC{以斜行冪內減二行差冪，餘為實。以二行差為從。一步常法。開平方得圓徑。}

\caoyueC{立天元一為圓徑。以斜行為弦，二行差為較。以斜行冪五萬七千六百步內減差冪六千四百步，餘五萬一千二百步為實。以差八十步為從。一步常法。開平方得一百二十步，即圓徑也。合問。}

\switchcolumn

\wenM{假設有一座圓形城池，甲、乙兩人同立於圓心。甲向東門直行，乙向南門直行。甲出東門後望見乙，便斜行與乙會合。只知甲直行二百步，乙直行比甲少八十步，又知斜行比甲直行多四十步。問圓城直徑多少？}

\daM{圓城直徑為一百二十步。}

\fayueM{以斜行步數的平方減去兩行步數之差的平方，餘數作為被開方數（實）。以兩行步數之差作為一次項係數（從）。以一步為二次項常數（常法）。開平方即得圓城直徑。}

\caoyueM{設天元一為圓城直徑。斜行為弦，兩行差為較。以斜行冪五萬七千六百步減去差冪六千四百步，餘五萬一千二百步，作為實。以差八十步作為從。以一步為常法。開平方得一百二十步，即圓城直徑。符合題意。}

\end{paracol}

\vspace{0.3cm}
\longline

%% ========================================================================
%% Problem 167 (Vol.12, Problem 17)
%% ========================================================================
\begin{paracol}{2}

\wenC{假令圓城一座，甲出南門直行，乙出東門直行。甲回望見乙，乃斜行與乙相會。只云甲行二百四十步，乙行八十步不及斜行一百六十步。問圓徑幾何？}

\daC{圓徑一百六十步。}

\fayueC{以甲行乘乙行，倍之為實。以甲行加乙行為從。一步常法。開平方得圓徑。}

\caoyueC{立天元一為圓徑。以甲行為股，乙行為勾。以勾股相乘得一萬九千二百步，倍之得三萬八千四百步為實。以甲行加乙行得三百二十步為從。一步常法。開平方得一百六十步，即圓徑也。合問。}

\switchcolumn

\wenM{假設有一座圓形城池，甲從南門直向行，乙從東門直向行。甲回頭望見乙，便斜行與乙會合。只知甲行二百四十步，乙行八十步，乙比斜行少一百六十步（即斜行比乙直行多一百六十步）。問圓城直徑多少？}

\daM{圓城直徑為一百六十步。}

\fayueM{以甲行步數乘乙行步數，加倍作為被開方數（實）。以甲行步數加乙行步數作為一次項係數（從）。以一步為二次項常數（常法）。開平方即得圓城直徑。}

\caoyueM{設天元一為圓城直徑。甲行為股，乙行為勾。以勾股相乘得一萬九千二百步，加倍得三萬八千四百步，作為實。以甲行加乙行得三百二十步，作為從。以一步為常法。開平方得一百六十步，即圓城直徑。符合題意。}

\end{paracol}

\vspace{0.3cm}
\longline

%% ========================================================================
%% Problem 168 (Vol.12, Problem 18)
%% ========================================================================
\begin{paracol}{2}

\wenC{假令圓城一座，甲乙二人同立於圓心。甲向東門直行，乙向南門直行。甲出東門望見乙，斜行與乙相會。只云甲直行一百八十步，乙直行一百二十步不及斜行六十步。問圓徑幾何？}

\daC{圓徑一百二十步。}

\fayueC{以甲行乘乙行，倍之為實。以甲行加乙行為從。一步常法。開平方得圓徑。}

\caoyueC{立天元一為圓徑。以甲行為股，乙行為勾。以勾股相乘得二萬一千六百步，倍之得四萬三千二百步為實。以甲行加乙行得三百步為從。一步常法。開平方得一百二十步，即圓徑也。合問。}

\switchcolumn

\wenM{假設有一座圓形城池，甲、乙兩人同立於圓心。甲向東門直行，乙向南門直行。甲出東門後望見乙，便斜行與乙會合。只知甲直行一百八十步，乙直行一百二十步，乙比斜行少六十步（即斜行比乙直行多六十步）。問圓城直徑多少？}

\daM{圓城直徑為一百二十步。}

\fayueM{以甲行步數乘乙行步數，加倍作為被開方數（實）。以甲行步數加乙行步數作為一次項係數（從）。以一步為二次項常數（常法）。開平方即得圓城直徑。}

\caoyueM{設天元一為圓城直徑。甲行為股，乙行為勾。以勾股相乘得二萬一千六百步，加倍得四萬三千二百步，作為實。以甲行加乙行得三百步，作為從。以一步為常法。開平方得一百二十步，即圓城直徑。符合題意。}

\end{paracol}

\vspace{0.3cm}
\longline

%% ========================================================================
%% Problem 169 (Vol.12, Problem 19)
%% ========================================================================
\begin{paracol}{2}

\wenC{假令圓城一座，甲出東門直行，乙出南門直行，丙出西門直行，丁出北門直行。只云甲行一百八十步，乙行一百二十步，丙行八十步，丁行四十步不及圓徑二十步。問圓徑幾何？}

\daC{圓徑一百六十步。}

\fayueC{以甲行乘乙行為實。以甲行加乙行為從。一步常法。開平方得圓徑。}

\caoyueC{立天元一為圓徑。以甲行為股，乙行為勾。以勾股相乘得二萬一千六百步為實。以甲行加乙行得三百步為從。一步常法。開平方得一百六十步，即圓徑也。合問。}

\switchcolumn

\wenM{假設有一座圓形城池，甲從東門直向行，乙從南門直向行，丙從西門直向行，丁從北門直向行。只知甲行一百八十步，乙行一百二十步，丙行八十步，丁行四十步且丁比圓徑少二十步（即圓徑比丁行多二十步）。問圓城直徑多少？}

\daM{圓城直徑為一百六十步。}

\fayueM{以甲行步數乘乙行步數，作為被開方數（實）。以甲行步數加乙行步數作為一次項係數（從）。以一步為二次項常數（常法）。開平方即得圓城直徑。}

\caoyueM{設天元一為圓城直徑。甲行為股，乙行為勾。以勾股相乘得二萬一千六百步，作為實。以甲行加乙行得三百步，作為從。以一步為常法。開平方得一百六十步，即圓城直徑。符合題意。}

\end{paracol}

\vspace{0.3cm}
\longline

%% ========================================================================
%% Problem 170 (Vol.12, Problem 20) —— FINAL PROBLEM
%% ========================================================================
\begin{paracol}{2}

\wenC{假令圓城一座，甲乙二人同立於圓心。甲向東門直行，乙向南門直行。甲出東門望見乙，斜行與乙相會。只云甲直行一百二十步，乙直行八十步，又云斜行與圓徑等。問圓徑幾何？}

\daC{圓徑八十步。}

\fayueC{以勾股冪相加為弦冪。以斜行為弦。以勾股相乘，倍之開平方得圓徑。}

\caoyueC{立天元一為圓徑即為弦。以甲行為股，乙行為勾。以勾冪六千四百步加股冪一萬四千四百步，共二萬八百步為弦冪。開平方得弦一百二十步。又以勾股相乘得九千六百步，倍之一萬九千二百步。以弦除之得一百六十步為實。一步常法。開平方得八十步，即圓徑也。合問。}

\switchcolumn

\wenM{假設有一座圓形城池，甲、乙兩人同立於圓心。甲向東門直行，乙向南門直行。甲出東門後望見乙，便斜行與乙會合。只知甲直行一百二十步，乙直行八十步，又知斜行與圓城直徑相等。問圓城直徑多少？}

\daM{圓城直徑為八十步。}

\fayueM{將勾的平方加上股的平方得到弦的平方。以斜行為弦。將勾股相乘，加倍後再開平方，即得圓城直徑。}

\caoyueM{設天元一為圓城直徑，也就是弦。甲行為股，乙行為勾。以勾冪六千四百步加上股冪一萬四千四百步，共二萬零八百步，作為弦冪。開平方得弦一百二十步。又以勾股相乘得九千六百步，加倍得一萬九千二百步。以弦除之得一百六十步，作為實。以一步為常法。開平方得八十步，即圓城直徑。符合題意。}

\end{paracol}

\vspace{0.3cm}
\longline


%% ========================================================================
%% Afterword
%% ========================================================================
\newpage
\section{後序 (Afterword) / 跋文}
\label{sec:afterword}

\longline

\begin{paracol}{2}
\leftlabel

\switchcolumn
\rightlabel

\end{paracol}

\vspace{0.3cm}
\longline

\begin{paracol}{2}

敬齋先生病且革，語其子克修曰：「吾平生著述，死後可盡燔去。獨測圓海鏡一書，雖九九小數，吾常精思致力焉，後世必有知者。庶可布廣垂永乎？」

\vspace{0.3cm}

先生於六藝百家，靡不貫串。文集近數百卷，常謙謙不自伐。惟於此書不忘稱異。予易簀之間，想有玄妙內得於心者。予以先生與先人同榜之故，素常兄事克修。克修兄命予重為序。予不敢說論艷藻，刻畫無鹽，唐突西子。直以所聞語之子弟，使知測圓海鏡之淵源云。

\switchcolumn

敬齋先生（李冶）病重臨終之際，對他的兒子克修說：「我平生的各種著述，死後可以全部燒掉。只有《測圓海鏡》這一部書，雖然只是關於算術的小技，但我曾經精心鑽研、傾注心血，後世必定有能理解它的人。但願它能夠廣為流傳、永久傳世吧？」

\vspace{0.3cm}

敬齋先生對於六藝百家之學，無不融會貫通。他的文集將近數百卷，但他始終謙遜，從不誇耀自己。唯獨對於這部書，他不忘稱讚其非凡之處。想來他在臨終之際，心中必定領悟到了某種深奧的玄妙。我因為先生與我的父親是同榜進士的緣故，一向像對待兄長一樣尊敬克修。克修兄命我重新寫一篇序。我不敢用華麗的辭藻來讚美，那樣做就像是給無鹽（醜女）畫上漂亮的妝容，反而是對西施（美女）的冒犯。我只是把所聽到的直接告訴後輩，讓他們知道《測圓海鏡》的來龍去脈罷了。

\end{paracol}

\vspace{1cm}
\longline

\vfill

\begin{center}
\fbox{\parbox{0.8\textwidth}{\centering
\textbf{\Large 測圓海鏡 卷十至卷十二 完}\\[0.5em]
{全書一百七十問終}\\[0.3em]
\textit{--- End of the Cèyuán Hǎijìng (Sea Mirror of Circle Measurements) ---}\\[0.3em]
{\small 文言文/白話文 雙語對照版 / Classical--Modern Chinese Bilingual Edition}
}}
\end{center}

\end{document}
%% ========================================================================
%% END OF DOCUMENT
%% ========================================================================
